\chapter{4x10\textsuperscript{9}D.C.}\label{4x10^9-6}

\hfill\break

Arrivò Diane seguita da En, una ventina di suoi cugini e altrettanti
estranei, tutti diretti verso il nuovo mondo. Jala fece entrare il
branco, quindi chiuse la porta scorrevole di acciaio ondulato del
magazzino. Le luci si affievolirono. Diane mi mise un braccio attorno al
collo e si lasciò accompagnare verso uno spazio relativamente pulito
sotto una delle lampade ad alogenuri. Ibu Ina srotolò un sacco di iuta
vuoto per farvela sdraiare sopra.

«Il rumore» sussurrò Ina.

Non appena si trovò in posizione orizzontale, Diane chiuse gli occhi,
sveglia ma visibilmente sfinita. Le sbottonai la camicetta e cominciai a
staccarla dalla ferita con delicatezza.

Dissi: «La mia borsa da medico\ldots»

«Sì, naturalmente.» Ina chiamò En e lo mandò di sopra a prendere
entrambe le borse, la mia e la sua. «Il rumore\ldots»

Quando iniziai a scollare il tessuto ruvido dal sangue incrostato della
ferita, Diane sussultò, ma non volevo medicarla prima di avere
controllato l'entità della lesione. «Quale rumore?»

«Appunto!» ribatté Ina. «Le banchine dovrebbero essere chiassose a
quest'ora del mattino. Invece c'è silenzio. \emph{Non c'è} rumore.»

Alzai la testa. Aveva ragione. Niente rumore, tranne il chiacchiericcio
nervoso dei Minangkabau e un tamburellare lontano dovuto al suono della
pioggia sull'alto tetto di metallo.

Ma non era il momento di preoccuparsene. «Vai a chiedere a Jala» la
esortai. «Scopri cosa sta succedendo.»

Quindi tornai a occuparmi di Diane.

ooo

«È superficiale» mormorò Diane. Inspirò profondamente. Stringeva gli
occhi per il dolore. «Almeno, \emph{credo} che sia superficiale.»

«Sembra un foro di proiettile.»

«Sì. La Reformasi ha scovato la casa sicura di Jala a Padang.
Fortunatamente eravamo sul punto di andarcene. \emph{Uh}!»

Aveva ragione. La ferita di per sé era superficiale, sebbene avesse
bisogno di una sutura. Il proiettile aveva attraversato il tessuto
adiposo poco sopra l'osso iliaco. Ma l'impatto, dove la pelle non era
lacerata, le aveva provocato delle brutte contusioni. Temevo che
l'ematoma fosse profondo e che il trauma avesse provocato danni interni.
Tuttavia, da quanto mi aveva detto, non aveva sangue nelle urine e,
considerate le circostanze, sia la pressione sanguigna che il polso
erano accettabili.

«Voglio darti qualcosa per il dolore, e dobbiamo suturare.»

«Se proprio devi, mettimi i punti ma non voglio narcotici. Dobbiamo
andarcene.»

«Non vorrai che ti suturi senza un anestetico\ldots»

«Allora anestesia locale.»

«Non è una clinica, non ho anestetici locali.»

«Allora cuci e basta, Tyler. Posso sopportare il dolore.»

Bene, ma io sarei stato in grado? Mi guardai le mani. Pulite, sì - nel
bagno del magazzino c'era acqua corrente e Ina mi aveva aiutato a
infilare i guanti di lattice prima di assistere Diane. Ero pulito e
qualificato. Ma per niente calmo.

Non sono mai stato schizzinoso sul lavoro. Anche quando ero studente di
medicina, persino durante le dissezioni, ero riuscito a tenere a freno
la spirale di empatia che ci fa sentire il dolore degli altri come se
fosse nostro. A fingere che l'arteria recisa che esigeva la mia
attenzione non fosse collegata a un essere umano. Fingere e crederci
veramente durante i pochi minuti necessari.

Però la mano aveva preso a tremarmi e l'idea di far passare un ago tra
quei lembi di pelle sanguinante mi pareva brutale, oltremodo crudele.

Diane mi appoggiò la mano sul polso per tenerlo fermo. «È roba da
Quarti» mormorò.

«Come?»

«Ti senti come se il proiettile avesse attraversato te al posto mio.
Vero?»

Annuii, sbalordito.

«È roba da Quarti. Credo che sia pensato per renderci persone migliori.
Ma rimani comunque un medico. Devi solamente concentrarti.»

«Se non ci riesco,» dissi, «lo faccio fare a Ina.»

Ma ci riuscii. In qualche modo, lo feci.

ooo

Ina tornò dopo essersi consultata con Jala. «Oggi è prevista una
manifestazione dei lavoratori» ci informò. «La polizia e la Reformasi
sono ai cancelli e hanno intenzione di prendere il controllo del porto.
Si prevede uno scontro.» Guardò Diane. «Come stai, mia cara?»

«In buone mani» bisbigliò Diane. La sua voce era sfinita.

Ina esaminò attentamente il mio lavoro. «Ben fatto» dichiarò.

«Grazie» risposi.

«Date le circostanze, intendo. Ma senti, ascoltami. Dobbiamo andarcene
molto in fretta. Ora come ora l'unica cosa che ci separa dalla prigione
è una rivolta dei lavoratori. Dobbiamo imbarcarci immediatamente sulla
\emph{Capetown Maru}.»

«Siamo ricercati dalla polizia?»

«Non credo cerchino voi, non nello specifico. Jakarta è entrata a far
parte di una sorta di accordo con gli americani per reprimere per quanto
possibile la tratta degli immigrati. Le banchine sono state rastrellate
in lungo e in largo in modo plateale, per fare una buona impressione su
quelli del consolato statunitense. Naturalmente non durerà a lungo. Ci
sono troppi soldi che passano di mano in mano perché la tratta venga
eliminata sul serio. Ma niente ha un effetto maquillage più dei
poliziotti in uniforme che trascinano via la gente dalla stiva delle
navi da carico.»

«Sono venuti nella casa sicura di Jala» disse Diane.

«Sì, sanno di te e del dottor Dupree, idealmente vorrebbero prendervi in
custodia, ma non è per questo che la polizia sta serrando i ranghi
presso i cancelli. Le navi stanno ancora lasciando il porto ma non
durerà a lungo. Il movimento sindacale è potente a Teluk Bayur. Hanno
intenzione di combattere.»

Jala era sulla porta d'ingresso e gridò qualcosa che non riuscii a
capire.

«Ora ce ne dobbiamo proprio andare» ci esortò Ina.

«Aiutami a fare una lettiga per Diane.»

Diane provò a tirarsi su. «Posso camminare.»

«No,» ribatté Ina, «su questo credo abbia ragione Tyler. Non provare a
muoverti.»

Raddoppiammo lo strato di iuta rattoppata in modo da formare una specie
di amaca su cui adagiarla. Presi un'estremità e Ina chiamò uno dei
Minangkabau più robusti perché afferrasse l'altra.

«Adesso sbrighiamoci!» gridò Jala, facendoci strada a gesti in mezzo
alla pioggia.

ooo

La stagione dei monsoni. Era un monsone? Il mattino sembrava il
crepuscolo. Nuvole come gomitoli di lana fradicia arrivavano dalle acque
grigie di Teluk Bayur, tagliando dalla vista le torri e i radar delle
enormi petroliere a doppio scafo. L'aria era calda e rancida. La pioggia
ci inzuppò mentre caricavamo Diane su un'automobile in attesa. Per il
suo gruppo d'immigrati Jala aveva organizzato un piccolo convoglio: tre
auto e un paio di piccoli carrelli merci.

La \emph{Capetown Maru} era ormeggiata all'estremità di un molo di
cemento a circa quattrocento metri di distanza. Lungo i pontili, in
direzione opposta, oltre le file di magazzini, depositi industriali e
grosse cisterne bianche e rosse dell'Avigas, una densa folla di
lavoratori portuali si era radunata presso i cancelli. Sotto la pioggia
tamburellante riuscivo a sentire qualcuno che strillava dentro un
megafono. Poi un suono che poteva essere, oppure no, quello
dell'esplosione di spari.

«Entra» mi esortò Jala, sollecitandomi a sedere sul sedile posteriore
dell'auto dove Diane era china sulle ferite come se stesse pregando.

«Andiamo, andiamo!» Salì al posto del guidatore.

Rivolsi un'ultima occhiata alla calca oscurata dalla pioggia. Un oggetto
grande come un pallone da football disegnò un'alta parabola sulla folla
lasciandosi dietro una scia di spire di fumo bianco. Un candelotto
lacrimogeno.

L'automobile fece uno scatto in avanti.

ooo

«C'è qualcun altro oltre alla polizia» ci informò Jala mentre
viaggiavamo lungo la lingua della banchina. «La polizia non sarebbe così
incosciente. Questa è la nuova Reformasi. Teppisti di strada in uniformi
governative, ingaggiati nei bassifondi di Jakarta.»

Uniformi e fucili. E altro gas lacrimogeno, nuvole torbide che si
confondevano con la foschia della pioggia. La folla iniziò a dipanarsi
ai margini.

Si udì un \emph{boom} lontano, e una palla di fuoco si innalzò in cielo
per alcuni metri.

Jala vide la scena dallo specchietto. «Mio Dio! Che idiozia! Qualcuno
deve avere fatto fuoco su un barile di petrolio. Le banchine\ldots»

Mentre costeggiavamo il molo, le sirene ululavano sull'acqua. Adesso la
folla era davvero in preda al panico. Per la prima volta riuscii a
vedere una fila di poliziotti farsi strada attraverso i cancelli
d'ingresso al porto. Le unità di punta imbracciavano armi pesanti e
indossavano maschere dal muso nero.

Un camion dei pompieri sbucò da un hangar diretto a sirene spiegate
verso il cancello.

Procedemmo verso una serie di rampe e ci fermammo nel punto in cui il
molo si innalzava verso il ponte principale della Capetown Maru.

La \emph{Capetown Maru} era un vecchio mercantile dipinto di bianco e
arancio ruggine che batteva bandiera di comodo. Tra il ponte principale
e il molo avevano appoggiato una corta passerella di alluminio e alcuni
dei primi Minangkabau la stavano già attraversando precipitosamente.

Jala saltò fuori dall'auto. Il tempo di portare Diane sul molo - non
avevamo più la lettiga di iuta e procedevamo a piedi, con lei appoggiata
pesantemente a me - e si era già immerso in un'accesa discussione in
inglese con l'uomo a capo della passerella: forse non proprio il
capitano della nave o il pilota ma comunque qualcuno che sembrava avere
un'autorità simile, un uomo tarchiato con il copricapo sikh e la
mascella serrata in modo cupo.

«Eravamo d'accordo mesi fa» sibilò Jala.

«Ma questo tempo\ldots»

«\emph{Indipendentemente} dal tempo!»

«Ma senza il permesso dell'Autorità Portuale\ldots»

«Sì, ma \emph{non c'è} nessuna Autorità Portuale\ldots{} guarda!»

Il gesto di Jala voleva essere retorico. Ma quando sventolò la mano in
direzione dei serbatoi di carburante vicino al cancello principale, una
delle cisterne esplose.

Non vidi nulla. L'impatto mi scaraventò sul cemento e percepii il calore
dell'esplosione sulla nuca. Il rimbombo fu fragoroso ma mi raggiunse in
un secondo momento. Appena fui in grado di muovermi mi rotolai sulla
schiena, le orecchie fischiavano. L'Avigas, pensai. O qualunque altro
materiale vi fosse contenuto. Benzene. Cherosene. Combustibile, perfino
olio di palma grezzo. Il fuoco doveva essersi esteso, oppure quei
poliziotti privi di addestramento avevano colpito il bersaglio
sbagliato. Mi girai cercando Diane e la vidi accanto a me che si voltava
indietro più disorientata che spaventata. Non riuscivo a sentire la
pioggia. Ma ci fu un altro suono, distintamente percettibile e davvero
minaccioso: il \emph{tintinnio} dei detriti in caduta. Schegge di
metallo, alcune infuocate. \emph{Tintinnavano}, mentre colpivano la
banchina di cemento o il ponte d'acciaio della \emph{Capetown Maru}.

«State giù,» urlava Jala, la sua voce mi giungeva sommersa, ovattata,
«state tutti giù, abbassate la testa!»

Provai a proteggere Diane facendole scudo con il mio corpo. Per pochi
interminabili secondi il metallo ardente ci cadde intorno come grandine
o finì nell'acqua nera oltre la nave. Poi semplicemente smise. Non scese
nient'altro che pioggia, lieve come il suono strusciato delle spazzole
sui piatti di una batteria.

Ci rialzammo. Jala stava già spingendo le persone sulla passerella
mentre guardava terrorizzato le fiamme. «Potrebbe non essere finita qui!
Salite a bordo tutti, andiamo, andiamo!» Dirigeva i compaesani oltre
l'equipaggio, che nel frattempo spegneva gli incendi sul ponte della
\emph{Capetown} e scioglieva le cime.

Il fumo si sollevò tutto intorno nascondendo le violenze perpetrate a
terra. Aiutai Diane a salire. A ogni passo sussultava e la ferita aveva
cominciato a trasudare dalle bende. Fummo gli ultimi a percorrere la
passerella. Un paio di marinai presero a ritirare la struttura
d'alluminio dietro di noi, mani sugli argani e sguardo alla colonna di
fuoco sulla riva.

I motori della \emph{Capetown Maru} rombavano sotto coperta. Jala mi
vide e si avvicinò per sollevare Diane dall'altro braccio. Lei si
accorse della sua presenza e chiese: «Siamo fuori pericolo?»

«No, lo saremo solo quando avremo lasciato il porto.»

Da una parte all'altra dell'acqua grigioverde risuonavano sirene e
fischi. Ogni nave in grado di partire si diresse in mare aperto verso
l'oceano. Jala si voltò verso il molo e si irrigidì. «I tuoi bagagli»
disse.

Erano stati caricati su uno dei piccoli carrelli merci. Due valigie
rigide, malandate e zeppe di documenti, medicinali e una scheda di
memoria. Giacevano ancora là, abbandonate.

«Fate tornare indietro quella passerella» ordinò Jala ai marinai del
ponte.

Sgranarono gli occhi dubitando della sua autorità. Il primo ufficiale
era andato sul ponte di comando. Jala gonfiò il petto e disse qualcosa
di tremendo in una lingua a me sconosciuta. I marinai scrollarono le
spalle e allungarono di nuovo la palanca verso il molo.

I motori della nave emisero un rombo più profondo.

Ripercorsi correndo la passerella, l'alluminio ondulato risuonava sotto
i piedi. Afferrai le valigie. Mi guardai indietro per l'ultima volta. A
terra, oltre il molo, un distaccamento di una decina di guardie della
Nuova Reformasi iniziò a correre verso la \emph{Capetown Maru}. «Mollare
gli ormeggi!» sbraitava Jala come se la nave fosse sua «Mollare gli
ormeggi, muoversi, presto!»

La palanca iniziò a ritirarsi. Lanciai il bagaglio a bordo e e subito
dopo vi saltai anch'io.

Raggiunsi il ponte prima che la nave cominciasse a muoversi.

Poi scoppiò un'altra cisterna della Avigas e fummo tutti scaraventati a
terra dall'impatto.